We construct a fixed regular model and an admissible weak-rank array whose causal probabilities stay separated although their complete observed two-sample laws merge.  This verifies the support, rank, span, and limiting conditions used in the contradiction.

Let
\[
\mathcal E=\{-1,+1\},\quad
\mathcal U=\mathcal W=\mathcal Y=\{0,1\},\quad
\mathcal X=\{x\}.
\]
Vectors indexed by a binary state space are ordered as \((1,0)\); the columns of \(S\), \(B_x\), and \(H_x\) are ordered as \((+1,-1)\), and the columns of \(M\) are ordered as \((1,0)\).  Define the fixed model \(\mathcal M^{(0)}\) by
\[
\pi=(1/2,1/2),\quad
S^{(0)}=
\begin{pmatrix}1/2&1/2\\1/2&1/2\end{pmatrix},\quad
M^{(0)}=
\begin{pmatrix}3/4&1/4\\1/4&3/4\end{pmatrix},
\]
\[
a_x(u)=1,\quad q^{(0)}=(1/2,1/2)^{\top},\quad
f^{(0)}_{x,1}(u,w)=f^{(0)}_{x,0}(u,w)=1/2.
\]
All entries of \(\pi,S^{(0)},M^{(0)},a_x,f^{(0)},q^{(0)}\) are strictly positive, every displayed kernel is stochastic, and \(\det M^{(0)}=1/2\ne0\).  The latent-shift factorization, target-mechanism invariance, strict primitive positivity, and proxy-channel injectivity therefore imply \(\mathcal M^{(0)}\in\mathfrak M_{+}\).  Marginalizing the factorization over \(U\) gives, for every \(e,w,y\),
\[
P_{O,0}(e,w,x,y)=\frac18,\quad
b_0=(1/2,1/2)^{\top},\quad
\theta_{x,1,0}=\frac12.
\]
Consequently,
\[
H_{x,0}=
\begin{pmatrix}1/4&1/4\\1/4&1/4\end{pmatrix},
\quad z_{x,1,0}=(1/4,1/4)^{\top}.
\]
The singular values of \(H_{x,0}\) are \(1/2\) and \(0\), so \(\operatorname{rank}(H_{x,0})=1\), \(\sigma_1(H_{x,0})=1/2\ge1/4\), and
\[
b_0=H_{x,0}(1,1)^{\top}.
\]
Thus the fixed-rank, retained-singular-value, and target-span conditions in \(\mathcal R_{1,1/4,1/16}\) hold.  Its source observed cells equal \(1/8\) and its target proxy cells equal \(1/2\), so the cell-floor condition holds with \(\eta=1/16\).

We also verify the positive-variance condition rather than assuming it.  At the baseline point,
\[
H_{x,0}^{\dagger}=
\begin{pmatrix}1&1\\1&1\end{pmatrix},
\quad
D_z\phi_1(H_{x,0},z_{x,1,0},b_0)
=(H_{x,0}^{\dagger}b_0)^{\top}=(1,1).
\]
Hence the source part of the scalar delta-method influence function has the form
\(\mathbf 1\{Y=1\}+\psi(E,W)\), up to a centering constant, where \(\psi\) is the contribution from the \(H\)-coordinates.  Under \(P_{O,0}\), conditional on \((E,W)\), \(Y\) is fair and independent of \((E,W)\).  The conditional-variance decomposition therefore gives
\[
\operatorname{Var}_{P_{O,0}}
 \bigl(\mathbf 1\{Y=1\}+\psi(E,W)\bigr)
\ge
\mathbb E_{P_{O,0}}
 \bigl[\operatorname{Var}_{P_{O,0}}(\mathbf 1\{Y=1\}\mid E,W)\bigr]
=\frac14.
\]
The source allocation multiplies this contribution by \(p^{-1}\), and the independent target contribution is nonnegative.  Therefore
\[
V_{1,x,1}(\mathcal M^{(0)};p)\ge\frac1{4p}>\frac1{16}.
\]
All defining conditions have been checked, so
\(\mathcal M^{(0)}\in\mathcal R_{1,1/4,1/16}\).

We now construct the admissible array.  Put \(\delta_n=\epsilon_n=n^{-2}\) and define \(\mathcal M_n^{(1)}\) using the same \(\pi\) and \(a_x\), but with
\[
S_n^{(1)}=
\begin{pmatrix}
1/2+\epsilon_n&1/2-\epsilon_n\\
1/2-\epsilon_n&1/2+\epsilon_n
\end{pmatrix},\quad
M_n^{(1)}=
\begin{pmatrix}
1/2+\delta_n&1/2-\delta_n\\
1/2-\delta_n&1/2+\delta_n
\end{pmatrix},
\]
\[
q^{(1)}=(3/4,1/4)^{\top},\quad
f^{(1)}_{x,1}(1,w)=3/4,\quad
f^{(1)}_{x,1}(0,w)=1/4,\quad
f^{(1)}_{x,0}=1-f^{(1)}_{x,1}.
\]
For \(n\ge2\), every primitive entry is strictly positive and
\[
\det M_n^{(1)}=2\delta_n>0,
\quad
\det S_n^{(1)}=2\epsilon_n>0.
\]
Thus each \(\mathcal M_n^{(1)}\) belongs to \(\mathfrak M_{+}\).  Since treatment is the singleton value \(x\), the posterior matrix equals \(S_n^{(1)}\).  Direct multiplication in the observable factorization yields
\[
B_{x,n}^{(1)}=M_n^{(1)}S_n^{(1)}=
\begin{pmatrix}
1/2+2\delta_n\epsilon_n&1/2-2\delta_n\epsilon_n\\
1/2-2\delta_n\epsilon_n&1/2+2\delta_n\epsilon_n
\end{pmatrix},
\]
\[
H_{x,n}^{(1)}=\frac12B_{x,n}^{(1)},\quad
\det B_{x,n}^{(1)}=4\delta_n\epsilon_n>0,
\quad
b_n^{(1)}=M_n^{(1)}q^{(1)}=
\begin{pmatrix}1/2+\delta_n/2\\1/2-\delta_n/2\end{pmatrix}.
\]
Invertibility of \(B_{x,n}^{(1)}\) implies
\(b_n^{(1)}\in\operatorname{col}(B_{x,n}^{(1)})\) in every row.  Under the fixed-proportion allocation and the row-wise independent source and target sampling assumptions, \(\mathbb M^{(1)}=(\mathcal M_n^{(1)})_{n\ge2}\) is therefore in \(\mathcal A_p\).  This is genuinely a weak-rank sequence: the unique solution of \(B_{x,n}^{(1)}\lambda_n=b_n^{(1)}\) satisfies
\[
\lambda_{n,1}+\lambda_{n,2}=1,
\quad
\lambda_{n,1}-\lambda_{n,2}=\frac1{4\epsilon_n},
\]
so its Euclidean norm, and the norm of the corresponding unconditional balancing vector, diverges.

Target-mechanism invariance and intervention at the singleton treatment give, in every row,
\[
\theta_{x,1,n}
=\sum_{u,w}q_u^{(1)}M_n^{(1)}(w,u)f^{(1)}_{x,1}(u,w)
=\frac34\frac34+\frac14\frac14
=\frac58.
\]
Thus the alternative causal probability is separated from the baseline value by \(1/8\).

It remains to compare the complete observed experiments.  For probability vectors on a common finite space, define
\(\lVert P-Q\rVert_{\mathrm{TV}}=\frac12\sum_a|P(a)-Q(a)|\).  Let \(S^{\star}\) and \(M^{\star}\) be the binary kernels all of whose entries are \(1/2\).  The auxiliary model using \(S^{\star},M^{\star}\), and \(f^{(1)}\) has observed marginal exactly \(P_{O,0}\): \(E\) and \(W\) are independent fair variables, and averaging \((3/4,1/4)\) over the fair latent variable makes \(Y\) fair and independent of \((E,W)\).  At the full-data level, adding and subtracting the density using \(S^{\star}\) and \(M_n^{(1)}\) gives
\[
\begin{aligned}
&\pi_eS_n^{(1)}(u,e)M_n^{(1)}(w,u)f^{(1)}_{x,y}(u,w)
-\pi_eS^{\star}(u,e)M^{\star}(w,u)f^{(1)}_{x,y}(u,w)\\
&\quad=\pi_e\bigl(S_n^{(1)}(u,e)-S^{\star}(u,e)\bigr)
M_n^{(1)}(w,u)f^{(1)}_{x,y}(u,w)\\
&\quad\quad+\pi_eS^{\star}(u,e)
\bigl(M_n^{(1)}(w,u)-M^{\star}(w,u)\bigr)f^{(1)}_{x,y}(u,w).
\end{aligned}
\]
Taking absolute values, summing all cells, and using stochasticity of the unchanged kernels bounds the full-data total variation by \(\epsilon_n+\delta_n\), because the columnwise total-variation distances of the two binary kernels are \(\epsilon_n\) and \(\delta_n\).  Marginalization cannot increase this absolute-sum bound, so
\[
\lVert P_{O,n}^{(1)}-P_{O,0}\rVert_{\mathrm{TV}}
\le\epsilon_n+\delta_n=2n^{-2}.
\]
The displayed target proxy vectors likewise give
\[
\lVert b_n^{(1)}-b_0\rVert_{\mathrm{TV}}=\delta_n/2.
\]
For probability vectors \(P_i,Q_i\), inserting one tensor factor at a time gives
\[
\bigotimes_{i=1}^N P_i-\bigotimes_{i=1}^N Q_i
=\sum_{j=1}^N
\left(\bigotimes_{i<j}P_i\right)\otimes(P_j-Q_j)\otimes
\left(\bigotimes_{i>j}Q_i\right).
\]
The triangle inequality and multiplicativity of the \(\ell_1\)-norm under tensor products imply that product total variation is at most the sum of the coordinatewise total variations.  Hence, using \(n_s(n)+n_t(n)=n\),
\[
\begin{aligned}
&\left\lVert
(P_{O,n}^{(1)})^{\otimes n_s(n)}\otimes(b_n^{(1)})^{\otimes n_t(n)}
-P_{O,0}^{\otimes n_s(n)}\otimes b_0^{\otimes n_t(n)}
\right\rVert_{\mathrm{TV}}\\
&\quad\le n_s(n)(\epsilon_n+\delta_n)+n_t(n)\delta_n/2
\le 2n^{-1}\longrightarrow0.
\end{aligned}
\]
This proves the asserted asymptotic indistinguishability of the entire supplied source and target samples.

Suppose, for contradiction, that a sequence of measurable nonempty confidence sets satisfies both requirements in the theorem.  Uniform coverage applied to the single admissible array \(\mathbb M^{(1)}\) gives
\[
\liminf_{n\to\infty}
\mathbb P_{\mathcal M_n^{(1)}}^{\,n_s(n),n_t(n)}
\{5/8\in\widehat C^{\mathrm{SAR}}_{x,1,n,1-\alpha}\}
\ge1-\alpha.
\]
Indeed, at each \(n\) the displayed probability is at least the infimum over \(\mathcal A_p\).  Total-variation convergence of the complete two-sample laws therefore transfers this event to the fixed baseline experiment:
\[
\liminf_{n\to\infty}
\mathbb P_{\mathcal M^{(0)}}^{\,n_s(n),n_t(n)}
\{5/8\in\widehat C^{\mathrm{SAR}}_{x,1,n,1-\alpha}\}
\ge1-\alpha>0. \tag{1}
\]

We next localize the regular Wald interval under \(\mathcal M^{(0)}\).  Every coordinate of \(\widehat H_{x,n}\), \(\widehat z_{x,1,n}\), and \(\widehat b_n\) is an empirical mean of a binary variable.  For an empirical mean based on \(N\) independent observations,
\[
\mathbb P\{|\widehat\mu-\mu|>t\}
\le t^{-2}\mathbb E[(\widehat\mu-\mu)^2]
\le\frac1{4Nt^2}.
\]
There are finitely many coordinates, and the allocation condition implies \(n_s(n)/n\to p\in(0,1)\) and \(n_t(n)/n\to1-p\in(0,1)\).  A finite union bound consequently gives joint \(O_P(n^{-1/2})\) convergence to \((H_{x,0},z_{x,1,0},b_0)\).  At that point the retained singular value is \(1/2\) and the next is \(0\).  Thus a fixed neighborhood lies in the strict-cutoff-gap region, where the selector-independent map \(\phi_1\), its derivative, and the finite-multinomial plug-in covariance map are continuous and bounded.  Since
\[
\phi_1(H_{x,0},z_{x,1,0},b_0)
=z_{x,1,0}^{\top}H_{x,0}^{\dagger}b_0=\frac12,
\]
the definitions of the Wald estimator and its totalized plug-in variance imply
\[
\widehat\theta^{\,W}_{x,1,n}=\frac12+O_P(n^{-1/2}),
\quad
\widehat V_n=O_P(1).
\]
The exceptional event outside that neighborhood has probability tending to zero, so the fixed zero fallback does not alter these stochastic orders.  After intersection with \([0,1]\),
\[
\Delta_n:=\sup_{\zeta\in I^W_{n,1-\alpha}}|\zeta-1/2|
=O_P(n^{-1/2}). \tag{2}
\]

Because \(\mathcal M^{(0)}\in\mathcal R_{1,1/4,1/16}\), the asserted regular-branch equivalence gives
\[
\sqrt n\,d_H\!\left(
\widehat C^{\mathrm{SAR}}_{x,1,n,1-\alpha},I^W_{n,1-\alpha}
\right)\longrightarrow_P0,
\]
and therefore the unscaled Hausdorff distance is \(o_P(1)\).  On the event
\(5/8\in\widehat C^{\mathrm{SAR}}_{x,1,n,1-\alpha}\), nonemptiness and the definition of \(d_H\) imply
\[
\frac18
\le d_H\!\left(
\widehat C^{\mathrm{SAR}}_{x,1,n,1-\alpha},I^W_{n,1-\alpha}
\right)+\Delta_n.
\]
By equation \((2)\) and the Hausdorff convergence, the right-hand side converges to zero in probability.  It follows that
\[
\mathbb P_{\mathcal M^{(0)}}^{\,n_s(n),n_t(n)}
\{5/8\in\widehat C^{\mathrm{SAR}}_{x,1,n,1-\alpha}\}
\longrightarrow0,
\]
contradicting equation \((1)\).  Therefore no measurable sequence of nonempty confidence sets, studentized or otherwise, can satisfy both uniform coverage over \(\mathcal A_p\) and first-order Hausdorff equivalence to the regular Wald interval on every fixed strongly identified submodel.  Any additional root-\(n\) diameter demand cannot repair this contradiction.

Finally, the stated positive half follows from the already established finite-sample concentration projection theorem.  For each row of every \(\mathbb M\in\mathcal A_p\), the defining array target-span condition holds.  Applying that theorem to \(\mathcal M_n\), \(n_s(n)\), and \(n_t(n)\) gives coverage at least \(1-\alpha\) in every row for the row-indexed concentration set.  Taking first the infimum over \(\mathcal A_p\) and then the lower limit preserves this bound.  Thus uniform coverage remains available, but the preceding indistinguishable pair proves that Wald equivalence requires a narrower array class, such as one imposing singular-value separation, bounded balancing weights, or comparable regularity.